% Part: many-valued-logic
% Chapter: syntax-and-semantics
% Section: matrices

\documentclass[../../../include/open-logic-section]{subfiles}

\begin{document}

\olfileid[id]{mvl}{syn}{mat}

\olsection{Matriks}

Suatu logika bernilai banyak didefinisikan oleh bahasanya, himpunan nilai
kebenarannya~$V$, suatu himpunan bagian yang beranggotakan nilai-nilai
kebenaran yang ditetapkan, dan fungsi kebenaran bagi setiap konektifnya.
Secara bersama-sama, unsur-unsur ini disebut suatu \emph{matriks}.

\begin{defn}[Matriks]
\ollabel{defn:matrix}
Suatu \emph{matriks} bagi logika~$\Log L$ terdiri atas:
\begin{enumerate}
\item suatu himpunan konektif yang membentuk bahasa~$\Lang L$;
\item suatu himpunan tak kosong~$V\neq\emptyset$ yang beranggotakan nilai
kebenaran;
\item suatu himpunan~$V^+\subseteq V$ yang beranggotakan nilai-nilai
kebenaran yang ditetapkan;
\item bagi setiap konektif $n$-tempat~$\star$ dalam~$\Lang L$, suatu fungsi
kebenaran~$\tf{\star}:V^n\to V$. Jika $n=0$, maka $\tf{\star}$ hanyalah
suatu unsur dari~$V$.
\end{enumerate}
\end{defn}

\begin{ex}
Matriks bagi logika klasik~\LogCL{} terdiri atas:
\begin{enumerate}
  \item Bahasa proposisional baku~$\Lang L_0$ dengan $\lfalse$, $\lnot$,
  $\land$, $\lor$, $\lif$.
  \item Himpunan nilai kebenaran~$V=\{\True,\False\}$.
  \item $\True$ merupakan satu-satunya nilai yang ditetapkan, yakni
  $V^+=\{\True\}$.
  \item Bagi $\lfalse$, kita mempunyai $\tf{\lfalse}=\False$. Fungsi
  kebenaran lainnya diberikan oleh tabel kebenaran biasa (lihat
  \olref{fig:tf-CL}).
\end{enumerate}
\begin{figure}
  \begin{center}
      \begin{tabular}{c|c}
        $\tf{\lnot}$ & \\
        \hline
        $\True$ & $\False$ \\
        $\False$ & $\True$
      \end{tabular}
      \quad
      \begin{tabular}{c|cc}
        $\tf{\land}$ & $\True$ & $\False$ \\
        \hline
        $\True$ & $\True$ & $\False$ \\
        $\False$ & $\False$ & $\False$
      \end{tabular}
      \quad
      \begin{tabular}{c|cc}
        $\tf{\lor}$ & $\True$ & $\False$ \\
        \hline
        $\True$ & $\True$ & $\True$ \\
        $\False$ & $\True$ & $\False$
      \end{tabular}
      \quad
      \begin{tabular}{c|cc}
        $\tf{\lif}$ & $\True$ & $\False$ \\
        \hline
        $\True$ & $\True$ & $\False$ \\
        $\False$ & $\True$ & $\True$
      \end{tabular}
    \end{center}
    \caption{Fungsi kebenaran bagi logika klasik~$\LogCL$.}
    \ollabel{fig:tf-CL}
  \end{figure}
\end{ex}

\end{document}
